\documentclass[11pt]{article}
\usepackage[T1]{fontenc}
\usepackage[margin=1in]{geometry}
\usepackage{amsmath,amssymb,mathrsfs,mathtools,booktabs,array,longtable}
\usepackage{microtype}
\usepackage[hidelinks]{hyperref}
\allowdisplaybreaks
\numberwithin{equation}{section}
\newcommand{\R}{\mathbb R}
\newcommand{\Z}{\mathbb Z}
\newcommand{\C}{\mathbb C}
\newcommand{\dd}{\mathrm d}
\newcommand{\Vol}{\operatorname{Vol}}
\newcommand{\tr}{\operatorname{tr}}
\newcommand{\Q}{\textup{[Q]}}
\newcommand{\AG}{\textup{[A]}}
\newcommand{\UR}{\textup{[R]}}
\newcommand{\US}{\textup{[S]}}
\newcommand{\qlabel}[1]{\nolinkurl{#1}}
\newcolumntype{P}[1]{>{\raggedright\arraybackslash}p{#1}}
\title{A digest of \texttt{QED3/qed3.tex}\\
and comparison with Dimock's lattice QED analysis}
\author{}
\date{22 October 2026}

\begin{document}
\maketitle

This note compares the local manuscript \cite{Q} with the supplied
versions of \emph{Covariant Axial Gauge} \cite{A},
\emph{Ultraviolet regularity for QED in $d=3$} \cite{R}, and
\emph{Ultraviolet stability for QED in $d=3$} \cite{S}.
These are abbreviated \Q, \AG, \UR, and \US.
References to their sections and equations use their own numbering;
selected searchable labels in \Q\ are collected at the end.
The comparison describes definitions, identities, and the stated proof
structure. It is not an independent verification of the long proofs in \Q.
In particular, ``the manuscript obtains'' below reports its claimed results.

\begin{center}
\begin{tabular}{@{}P{0.24\textwidth}P{0.70\textwidth}@{}}
\toprule
Question & Answer at the level of the actual definitions\\
\midrule
Same lattice model? & No. The local Maxwell term is shared, but the
fermion graph and operator, flux sectors, and treatment of harmonic
gauge fields differ.\\[5pt]
Same RG steps? & The same block-averaging and Gaussian-integration
framework, with particularly close photon formulas. The full maps
differ through the chamber first step, external-source bookkeeping,
cutoffs, and volume-comparison construction.\\[5pt]
Same small/large fields? & Closest to the global bounded-field analysis
of \UR\ in the small-field part. The large-field argument in \Q\ is
different from the spatial, multiscale expansion in \US.\\[5pt]
Same effective actions? & The same kinetic/mass/vacuum/polymer
decomposition in the small-field sector, but different operators,
activity spaces, normalizations, and source dependence. The full
large-field density of \US\ has additional history and boundary data.\\
\bottomrule
\end{tabular}
\end{center}

\section{What the manuscript is trying to construct}

The observable in \Q\ is the real field strength $F$, regarded in the
continuum as a tempered two-form on $\R^3$. For a fixed positive reference
mass $0<\bar m\le1$ and sufficiently small physical charge $e$, the
manuscript chooses a cutoff-dependent real mass shift $m_{0,N}(e)$.
Its regularization has
\begin{equation}
 \eta=L^{-N},\qquad q=e\sqrt\eta,\qquad
 \mu_{\mathrm{lat}}=\bar m\eta+m_{0,N}(e),\qquad
 R=8hL^T,\quad h=L^{\lfloor\log_L\bar m^{-1}\rfloor}.
\end{equation}
It takes the infinite-volume limit at fixed $\eta$ first, then takes
subsequential limits as $\eta\downarrow0$. The main theorem concerns
these limit laws and their translation-ergodic components; it does not
assert uniqueness of the limit.

The stated conclusion is that almost every such component satisfies
the tensor Osterwalder--Schrader axioms, including clustering, and is
non-Gaussian. The parent law is asserted to satisfy the other axioms,
but clustering is asserted for the components. Two quantitative inputs
are Gaussian domination and a positive fourth cumulant:
\begin{align}
 \int e^{F(g)}\,\dd\mu_\xi(F)&\le e^{\|g\|_2^2/2},\\
 \kappa_4^{\mu_\xi}(f)&=e^4b(f)+O_f(e^{4+\sigma}),\qquad b(f)>0,
 \label{eq:digest-cumulant}\\
 b(f)&=-6\tr\big[(S_{\bar m}\,i\gamma\cdot\dd^*\beta)^4\big],
 \quad f=\dd\dd^*\beta,
 \quad S_{\bar m}=(\bar m+\gamma\cdot\partial)^{-1}.
\end{align}
The charge threshold for positivity may depend on the chosen probe
$\beta$. This is a statement about field-strength Schwinger functions,
not a construction here of every charged fermion or gauge-potential
correlation function.

There are two substantial parts to the argument.
First, chamber geometry supplies reflection-positive crossings.
Slab traces, fold inequalities, and comparison with adapted tori are
used to obtain positivity at the three coordinate and six diagonal
mirror families in the volume limit (\Q, Sections 2--4).
Gaussian domination gives tightness and tensor moment bounds
(Sections 5--6). The manuscript then uses reflection-positive
reconstruction, a sharp spectral cone, analytic continuation and
angular-frequency arguments, and the Bianchi identity to obtain full
$O(3)$ covariance and componentwise clustering (Sections 7--8).
It explicitly does not infer full rotational covariance just from the
finite cubic reflection group.

Second, the block integration in Sections 9--16 tunes the relevant
couplings and controls polymer activities. Sections 17--23 compute
free and interacting external-field responses on reference tori.
Sections 24--26 choose compatible constants and transfer reference
estimates, including products of translated Wick polynomials, to large
volumes and then to ergodic components. A positive cumulant for the
parent mixture alone would not suffice for the componentwise claim.
Sections 27--28 assemble the tensor axioms and
\eqref{eq:digest-cumulant}.

The three Dimock papers have narrower and distinct roles.
\AG\ develops the photon gauge fixing and block-integration machinery;
it is not a fermionic QED construction. \UR\ controls a globally
bounded-field RG flow and tunes vacuum energy and mass to prescribed
terminal values (Theorems 1 and 3). \US\ incorporates the large-field
corrections and proves, for the specified counterterms and small charge,
\begin{equation}
 \frac12\le\frac{Z(N,e)}{Z(N,0)}\le\frac32
 \qquad\text{uniformly in }N
 \label{eq:digest-stability}
\end{equation}
in fixed finite physical volume (\US, Theorem 2, (42)). That stability theorem alone is not
the volume-first, non-Gaussian field-strength OS theorem stated in \Q.

\section{The lattice models are different}

\subsection{The common local structure and scaling}

Both constructions use a real abelian connection, a quadratic Maxwell
energy, four-component fermionic spin variables, and gauge-covariant
nearest-neighbor hopping on their respective fermion graphs. In
Dimock's physical lattice notation, with mesh $\eta$, the action is
\begin{equation}
 S_{\mathrm D}(A,\bar\psi,\psi)
 =\tfrac12\|\dd A\|^2
  +\langle\bar\psi,(D_A^{\mathrm W}+\bar m)\psi\rangle
  +m_N\langle\bar\psi,\psi\rangle+\varepsilon_N\Vol,
 \qquad
 D_A^{\mathrm W}=\gamma\cdot\nabla_A-\tfrac\eta2\Delta_A.
 \label{eq:digest-Wilson}
\end{equation}
The Wilson coefficient in this definition is one; covariant hopping
has phase $e^{ie\eta A(b)}$ (\UR, Section 1.2, (3)--(12)).
The starting physical torus in \UR\ and \US\ has side one.
After rescaling to unit mesh, both schemes have the same engineering
dimensions:
\begin{equation}
 e_k=eL^{(k-N)/2},\qquad
 \bar m_k=\bar m L^{k-N},\qquad
 A_L(b)=L^{-1/2}A(b/L),\qquad
 \Psi_L(x)=L^{-1}\Psi(x/L).
 \label{eq:digest-scaling}
\end{equation}
In \Q\ the reference mass is called $\mu_k$; its physical level length
is $a_k=L^{k-N}$. The inverse-power mesh $\eta_k=L^{-k}$ used for
reconstructed fine fields at level $k$ is a different quantity from
$a_k$ and from the original physical cutoff $\eta$.
The small parameter called $\varepsilon$ in activity bounds is also
different from the vacuum energy and from the lattice spacing.

\subsection{Chamber fermions are not Wilson fermions on the cubic sites}

In \UR\ there is one four-component spinor per cubic site. In \Q\ a
photon site $c\in\Z^3$ carries a cluster of $48$ chamber sites
$(c,g)$, $g\in O_{\mathrm{cub}}$. Each chamber site carries four spin
components, so the fine Dirac matrix has dimension $192V$, rather
than $4V$, on a photon torus of $V$ sites.
There are four chamber neighbors, with standard weights
\begin{equation}
 (w_0,w_1,w_2,w_3)=(6,6\sqrt2,6\sqrt2,6).
\end{equation}
Three types of bond stay within a cluster and have phase one; the
fourth projects to a cubic photon link. Eight chamber bonds project
to each photon link. The hopping matrices are
\begin{equation}
 w(x,x^{(i)})=\frac{w_i}{2}
       (\rho-\gamma\cdot\widehat d_i(x)),\qquad 0<\rho<1,
 \quad
 (\mathcal D_A\psi)(x)
 =M\psi(x)-\sum_i w(x,x^{(i)})U_A(x,x^{(i)})\psi(x^{(i)}),
\end{equation}
where $U_A(x,x')=e^{iqA(\pi x,\pi x')}$ and the mass after the diagonal
shift is $m=M-\frac\rho2\sum_iw_i$ (\Q, Sections 2.2--2.3).
This geometry is chosen so that a fermion bond crossing any of the
nine mirror families is perpendicular to that mirror and connects
reflected sites. The reflected crossing coefficients are
$w_i(1\pm\rho)/2$.

There is a precise Wilson comparison, but it does not identify the
operators. Let $J$ inject a four-spinor as a constant on its cluster,
let $J^*$ average over the $48$ labels, and let $Q_{\mathrm h}=1-JJ^*$.
For the free unit-mesh symbol, \Q, Section 9.2 gives
\begin{align}
 D_{\mathrm{ll}}(p)&=m+i\gamma\cdot\sin p+
                         \rho\sum_\mu(1-\cos p_\mu),\\
 B_m(p)&=D_{\mathrm{ll}}(p)
         -D_{\mathrm{lh}}(p)D_{\mathrm{hh}}(p)^{-1}D_{\mathrm{hl}}(p)
       =m+i\gamma\cdot\sin p+O(|p|^2).
 \label{eq:digest-Schur}
\end{align}
The compression $D_{\mathrm{ll}}$ is Wilson-like. The operator after
eliminating the heavy modes is the Schur complement $B_m$, and it
has the additional displayed term. The heavy modes also contribute
to the determinant. Thus neither taking a cluster average nor noting
the same leading Dirac symbol proves equality of the finite-cutoff
models. The manuscript's heavy-gap and propagator estimates are
additional input; they are not the cubic Wilson estimates with names
changed. Agreement of the leading symbol is also not a proof that
the two interacting regularizations have the same continuum law.

\subsection{Fluxes and harmonic modes}

In Dimock's model, $A$ is a globally defined real periodic one-cochain
and $F=\dd A$. Its flux through each coordinate two-torus is zero.
The harmonic one-form directions, which have zero Maxwell action,
require separate treatment from ordinary exact gauge directions.
\AG, Section 2.4, (21)--(22), and Section 3.1, (39)--(40), suppress
torons by imposing zero averaged circulations. \US, Section 4.1,
(409), likewise includes $\delta(Q^*A_K)$ in the final integral.
\AG\ explicitly distinguishes this suppression from gauge fixing.

The configuration space in \Q, Section 2.1 is instead
\begin{align}
 \operatorname{Conf}(\Gamma)
 &=\{(A,z):A\in C^1(\Gamma;\R),\ z\in C^2(\Gamma;\Z),\ \dd z=0\}/\sim,
 \\
 (A,z)&\sim(A+2\pi k/q,z-\dd k),\qquad k\in C^1(\Gamma;\Z),\\
 F&=\dd A+\frac{2\pi}{q}z,\qquad u_b=e^{iqA(b)}.
 \label{eq:digest-flux-model}
\end{align}
It sums over all closed integer $z$ and integrates $A(b)$ over
$[0,2\pi/q)$. The curvature is unbounded and closed. In particular,
the closedness constraint excludes monopole defects; the formula
should not be identified with a compact Wilson plaquette action or
with an unconstrained integer-plaquette sum.
Its finite-volume weight after integrating the fermions is
\begin{equation}
 e^{-\|F\|^2/2}\det\mathcal D_A\,\dd(A,z).
\end{equation}
The paired spin blocks give a nonnegative determinant for every real
mass, including when the tuned total lattice mass is negative
(\Q, Section 2.3).

Even restricting \Q\ to zero total flux does not make its global
photon prescription identical to Dimock's. The zero-flux sector
unfolds as
\begin{equation}
 Z^{(0)}[\Phi]
 =\mathcal N_n\int_{[-\pi,\pi]^3}\frac{\dd\theta}{(2\pi)^3}
  \int_{V_{\mathrm t}}\Phi(\dd A)e^{-\|\dd A\|^2/2}
        \det\mathcal D_{A+H_\theta}\,\dd\sigma(A).
 \label{eq:digest-unfold}
\end{equation}
Here $V_{\mathrm t}$ is the transverse, mean-zero potential space and
$H_\theta$ is a constant potential representing the three holonomies
(\Q, Section 20.3). Thus the manuscript integrates over torons, rather
than setting them to zero. The local Maxwell Gaussian on exact
curvatures is the common part. The fermion determinant can depend on
the holonomy, so this distinction is not just a field-independent
normalization.

\UR, Section 1.2 allows periodic or antiperiodic fermions. The main
cutoff tori in \Q\ use fermions antiperiodic in all three directions;
its estimates also track twists and other torus geometries. These
boundary choices should be matched explicitly in any comparison.

\section{Which parts of the RG step agree?}

\subsection{A common exact Gaussian framework}

Before adding cutoffs, Dimock's simultaneous block step has the form
\begin{equation}
 \widetilde\rho_{k+1}(A',\Psi')
 =\int\delta(A'-QA)\,\delta(\tau A)\,
       \delta_G\bigl(\Psi'-Q(\widetilde A)\Psi\bigr)
       \rho_k(A,\Psi)\,\mathrm D\Psi\,\dd A,
 \label{eq:digest-RG}
\end{equation}
followed by the rescaling in \eqref{eq:digest-scaling} and its measure
normalizations (\UR, (392)--(393); \US, (4)--(5)). The gauge average
is a straight-path average over an $L$-block, and $\tau$ averages the
six coordinate-order paths from the block center. The fermion block
factor is a Gaussian in Grassmann variables, with strength $b/L$,
not a sharp delta constraint. Its covariant average uses the
exponential of an averaged line integral, not the average of six
separate exponentials (\UR, (38)--(48)).

The photon definitions in \Q, Section 10 are closely aligned with
\AG\ and \UR. Its cochain averages commute with $\dd$ in the
appropriate degrees. The Landau background $H_k^{\mathrm{ph}}A$
minimizes the Maxwell energy subject to
\begin{equation}
 \mathbf Q_k\mathscr A=A,\qquad
 \mathcal R_k\dd^*\mathscr A=0,
 \quad \mathcal R_k\text{ projects onto }
           \Delta_{\mathrm{sc}}(\ker\mathbf Q_k).
\end{equation}
This is the constrained Landau prescription of \UR, (186)--(190).
The one-step fluctuation is in the axial slice
$\mathbf QZ=0$, $\tau Z=0$. Using a local coordinate map $C$,
\begin{equation}
 \Delta_k=(\dd H_k^{\mathrm{ph}})^T\dd H_k^{\mathrm{ph}},\qquad
 C_k=(C^T\Delta_k C)^{-1},\qquad
 Z=CC_k^{1/2}\zeta,\quad \zeta\sim\mu_I.
 \label{eq:digest-photon}
\end{equation}
The uncut gauge-invariant fluctuation integral therefore has the
same form as \AG, (152):
\begin{equation}
 F_{k+1}(\mathscr A)
 =\int F_k\bigl(\mathscr A_L+
                 H_k^{\mathrm{ph}}CC_k^{1/2}\zeta\bigr)\,\dd\mu_I(\zeta).
\end{equation}
The simultaneous use of a regular Landau background and positive
axial fluctuation coordinates is shared machinery, not a departure
introduced by \Q. The manuscript makes coordinate choices, local
regularity bounds, and compatibility on covers explicit. Its actual
nonlinear step also restricts this integral to a box, as discussed below.

The fermionic Schur-complement identity has the same algebraic form
in both analyses. With $\mathcal Q_k$ the applicable fine-to-coarse
average and $\mathcal Q_k^\sharp(A)=\mathcal Q_k(-A)^T$, write
\begin{align}
 S_k(A)&=(D_A+\mu_k+b_k\mathcal Q_k^\sharp\mathcal Q_k)^{-1},\\
 H_k(A)&=b_kS_k(A)\mathcal Q_k^\sharp(A),\\
 D_k(A)&=b_k-b_k^2\mathcal Q_k(A)S_k(A)\mathcal Q_k^\sharp(A).
 \label{eq:digest-fermion-block}
\end{align}
In \UR\ these are (78)--(84), with $D_A$ the cubic Wilson operator.
In \Q\ the underlying $D_A$ is the chamber operator and
$\mathcal Q_k=Q_kJ^*$. At level zero the manuscript uses the special
convention $D_0=D_A+\mu_0$, $H_0=1$; the first one-step average is
$Q(A)J^*$. That step integrates all $192n_0$ incoming pairs of
Grassmann generators into $4n_1$ outgoing pairs, including the heavy
components and their determinant. Later coarse fields have four
components per coarse site. See \Q, Sections 9.5, 11.5, and 14.7.

\subsection{The shared relevant flow, and the changes in its implementation}

Both schemes extract vacuum energy and a scalar fermion mass from
the polymer activity, and reblock and rescale the remainder. At zero
external field the recursions have the common structure
\begin{align}
 m_{k+1}&=L\bigl(m_k+m(E_k)\bigr),\\
 \varepsilon^{\mathrm v}_{k+1}
   &=L^3\bigl(\varepsilon^{\mathrm v}_k+
                  \varepsilon(E_k)+\delta\varepsilon_k^0\bigr),\\
 E_{k+1}&=\mathcal L(RE_k)^{\mathrm{adj}}
                   +\mathcal L E_k^{\mathrm{det}}+E_k^*(m_k,E_k).
 \label{eq:digest-flow}
\end{align}
Placement of $\mathcal L$ in the notation for the last two terms
varies between the sources. The adjustment preserves the evaluated
activity while improving its quadratic-kernel representation.
In \US\ the tiny vacuum correction is tracked separately from
the main vacuum coupling.
Compare \UR, (401)--(404), \US, (28), (33)--(34), and \Q, Section 14.7.
The same characteristic contraction factor
$O(L^{-1/4+2\varepsilon})$ appears in \UR\ and \Q.

The extraction is substantive agreement. For small polymers,
$d_M(X)\le L$, the constant coefficient at zero connection and the
integrated ordinary quadratic kernel are subtracted. Consequently
\begin{equation}
 (RE)_{00}(X;0)=0,\qquad
 \int_{X\times X}(RE)_{20}(X;0;x,y)\,\dd x\,\dd y=0.
\end{equation}
After summing over small polymers through a cube, symmetries make the
mass matrix scalar and real. These are the same normalization
conditions and relevant directions as in \UR, (317)--(326).
In \Q\ the integrals additionally average cluster labels with measure
$\eta_k^3/48$ and the extraction is evaluated at both connection slots
equal to zero (Section 12.3). Large polymers are left unsubtracted;
their geometric decay is used instead.

Neither scheme introduces a separately tuned charge or photon-mass
counterterm in this flow. The charge obeys
$e_{k+1}=L^{1/2}e_k$. This bookkeeping does not assert that the photon
effective action has no polarization corrections: field-dependent
determinants and activities retain such corrections.

The boundary-value construction for the relevant flow also agrees
in method: activities start at $E_0=0$, masses are solved backwards
from a terminal value, and the vacuum is obtained by a backwards
scalar sum. \UR\ chooses $m_K=\varepsilon_K=0$
(Section 7, Theorem 3); \US\ uses the corresponding controlled flow
with a stopping scale satisfying $L^{N-K}=Mr_K$.

The manuscript adds a specific comparison of local operations on
odd covers (Sections 15--16). It selects one bare pair on a prescribed
top torus of physical side
\begin{equation}
 R_{\mathrm{top}}=hMr_K,\qquad
 R^*=R_{\mathrm{top}}/L,\qquad K=N+\log_Lh.
\end{equation}
These physical sides are cutoff-independent for fixed charge and
reference mass. It then claims exact agreement of extracted
couplings, and of local activities on suitable lifted polymers,
on the smaller tori used in the construction. The terminal zeros are
imposed on the top torus; the couplings need not vanish when a smaller
torus stops. This shared tuning and the later transfer to arbitrarily
large cutoff volumes are additional arguments, not consequences of
the fixed-volume flow theorem in \UR.

\section{Small and large fields: the main methodological difference}

\subsection{Three different uses of ``small''}

It is useful to distinguish small \emph{field amplitudes}, the larger
complex \emph{analyticity domains}, and small \emph{polymers}.
The latter is the geometric condition $d_M(X)\le L$; a large polymer
can still be entirely in a small-field region.
Grassmann fields have coefficient norms rather than pointwise
probabilistic amplitude cutoffs.

Both \UR\ and \Q\ use logarithmic real-field thresholds
\begin{equation}
 p_k=(-\log e_k)^p,\qquad p_{0,k}=(-\log e_k)^{p_0},
 \qquad p_0<p,
\end{equation}
and complex connection domains of the form
\begin{equation}
 \|A\|_\infty<e_k^{-3/4+\varepsilon},\qquad
 \|\partial A\|_\infty<e_k^{-3/4+2\varepsilon},\qquad
 \|\delta_\alpha\partial A\|_\infty<e_k^{-3/4+3\varepsilon}.
 \label{eq:digest-domain}
\end{equation}
The extended domain $\widetilde R_k$ allows a real background which
has a small representative in a local gauge, plus a globally small
complex part (\UR, Definitions 1--2; \Q, Sections 10.8 and 12.2).
Thus ``small field'' is not the gauge-dependent assertion that every
real potential, including its harmonic part, is globally small.
Curvature bounds and local gauges connect the logarithmic cutoffs to
these analytic domains. \US\ uses spatially multiscale analogues
$R_{k,\Omega}$ and $\widetilde R_{k,\Omega}$, with additional margin
factors (Definitions 1--2).

\subsection{The bounded-field approximation in the regularity paper}

\UR\ replaces the exact integral \eqref{eq:digest-RG} by one with
$\chi_k\widehat\chi_k$ at each step. Its stated cutoffs are (396):
\begin{equation}
 \chi_k=\mathbf1\{\|\dd\mathcal A_k\|_\infty\le p_k\},\qquad
 \widehat\chi_k=
  \mathbf1\{\|A_k-H_k^xA_{k+1}\|_\infty\le p_{0,k}\},
 \quad \mathcal A_k=H_kA_k
 \label{eq:digest-D-cutoffs}
\end{equation}
on the whole torus. Its theorem controls this modified flow. The
large-field complement is not established by that theorem; restoring
it is a task of \US.

\subsection{The spatial large-field expansion in the stability paper}

\US\ retains all configurations by summing over histories
\begin{equation}
 \Pi=(\Pi_1,\ldots,\Pi_k),\qquad
 \Pi_j=(\Omega_j,\Lambda_j;P_j,Q_j,R_j,U_j),\qquad
 \Omega_1\supset\Lambda_1\supset\cdots\supset\Omega_k\supset\Lambda_k.
\end{equation}
The regions are assembled from scale-dependent $Mr_j$ cubes,
where $r_j$ is a rounded power of $-\log e_j$. The bad sets record
violations of the successive field and fluctuation restrictions;
buffer layers separate the active regions. Older-scale gauge and
fermion fields are retained outside the current active region.
The representation includes multiscale minimizers, unintegrated
fields, characteristic functions, and boundary activities $B_{k,\Pi}$
(\US, Section 2.1, (6)--(27)).

The bulk activity $E_k(X)$ is independent of the history when it lies
in the small-field region. Relevant extraction is confined there;
boundary corrections and large-field terms are controlled separately.
Small exponential factors from violated cutoffs pay for summing all
region histories. At the last scale the torus is a single $Mr_K$ cube:
the terminal region is either the whole torus or empty. The partition
function splits into the all-small contribution and the remaining
histories (\US, Section 4.1, (410)--(412)); the latter require estimates
of their fermionic as well as bosonic integrals (Section 4.2).

\subsection{The manuscript's global event and its complement}

The step used in \Q\ restricts the independent photon coordinates
in \eqref{eq:digest-photon} to
\begin{equation}
 \mathrm{Box}_j=\{\zeta:|\zeta(\Upsilon)|\le p_{0,j}
                              \text{ for every coordinate }\Upsilon\}.
\end{equation}
This is a cutoff on whitened Gaussian coordinates, rather than the
literal potential-fluctuation cutoff in \UR, (396). They serve
similar purposes but define different truncated integrals. The box
mass is retained as a tiny vacuum contribution:
\begin{equation}
 \mu_I(\mathrm{Box}_j)
   =e^{-\ell_j^3\delta\varepsilon_j^0},\qquad
 \delta\varepsilon_j^0\le12e^{-p_{0,j}^2/2}.
\end{equation}
Domain preservation propagates the terminal small-field condition
down through these boxes. The resulting event is
\begin{equation}
 \mathfrak G=\bigl\{\zeta_j\in\mathrm{Box}_j\ (j<k_*),\quad
                     \|\dd A_*\|_\infty\le p_*\bigr\}.
 \label{eq:digest-good-event}
\end{equation}
The exact iterated density in \Q, Section 22.2 represents the original
fine density restricted to $\mathfrak G$. It does not by itself
represent the whole measure after discarding the restriction.

To control $\mathfrak G^c$, the manuscript uses a separate Gaussian
domination result for the \emph{normalized, harmonic-averaged,
zero-flux interacting law} $\bar\mu$:
\begin{equation}
 \bar\mu(e^{\langle g,F\rangle})
  \le \exp\bigl(\|\operatorname{Pr}_{F_N}g\|^2/2\bigr),
 \qquad F_N=\dd C^1.
 \label{eq:digest-zero-domination}
\end{equation}
This is \Q, Section 21.5. The all-flux Gaussian bound does not
automatically imply a bound on this conditional law. The manuscript
retains a flux projection in its coordinate-slice trace to establish
\eqref{eq:digest-zero-domination}. Nor is the displayed statement a
uniform claim for every fixed twist separately.

The photon tower coordinates are linear functionals of curvature
with controlled row norms. Exponential Markov bounds and a union
bound then give (\Q, Section 22.3)
\begin{align}
 \bar\mu(\mathfrak G^c)
 &\le4\ell_*^3\sum_{i=1}^{k_*}L^{3i}
   \exp\left[-\tfrac12
          (v_*+\tfrac i2\log L)^{2p_0}\right]
       +6\ell_*^3e^{-p_*^2/2}\\
 &\le6\ell_*^3
       \left(e^{-p_{0,*}/2}+e^{-p_*^2/(2C_0^2)}\right),
 \qquad v_*=-\log e_*.
 \label{eq:digest-tail}
\end{align}
This argument does not require independence under the interacting
law. At fixed reference geometry, the scale sum absorbs the growing
number of ultraviolet coordinates, and $\ell_*,e_*$ are independent
of the ultraviolet cutoff after the specified threshold.
For external-response derivatives, Section 23.4 additionally uses
shifted events, real-source bounds, and holomorphic estimates; small
probability at zero source alone would not bound those derivatives.

The volume factor $\ell_*^3$ matters. This is not a direct argument
that all fields remain globally small with high probability on
arbitrarily large tori at fixed charge. The manuscript instead uses
the two reference tori $R^*$ and $LR^*$, establishes a common pressure,
and transfers local and product estimates by chessboard and partition
comparison inequalities (Sections 22.6 and 25).
Flux sectors are then handled as a separate part of that comparison:
\begin{equation}
 P_n(w)\le e^{-2\pi^2|w|^2/(q^2n)},\qquad q^2n=e^2R.
\end{equation}
This yields useful zero-sector dominance on the chosen reference
scales when $e^2R$ is small. It does not say nonzero flux is negligible
uniformly as $R\to\infty$. The transfer argument is what connects
the reference estimates to the all-flux volume-first construction.

Consequently \Q\ does not reproduce the history sum of \US.
Its proposed replacement depends on reflection positivity, the
sector-specific domination bound, and transfer between volumes.
The success of Dimock's large-field analysis does not independently
validate that replacement.

\section{Effective actions: a dictionary, not an identity}

\subsection{The meaning of the effective density}

Here ``effective action'' refers to the density obtained by integrating
short-distance variables while retaining coarse fields. These are
Wilsonian block densities, not one-particle-irreducible effective
actions defined by a Legendre transform. Three objects must be kept
distinct: the full blocked density $\rho_k$, its chosen quadratic
reference action, and its polymer correction $E_k$. In particular,
$E_k$ is not the whole effective action, and $\log\det D_A$ is not
the whole blocked density.

In \AG, (61)--(62), the photon density is
\begin{equation}
 \rho_k(A_k)=Z_k F_k(H_k^xA_k)
               e^{-\langle A_k,\Delta_kA_k\rangle/2}.
\end{equation}
The arbitrary interaction is retained in $F_k$; no fermionic mass
extraction or complete interacting QED action is defined by this
formula alone. Even the free blocked photon quadratic form is
$\Delta_k$, not simply the original unit-lattice $\dd^*\dd$.

\UR, (397)--(399), writes the global small-field density as
\begin{equation}
 \begin{aligned}
 \rho_k(A_k,\Psi_k)=\mathcal N_k\exp\Bigl\{
   &-\tfrac12\|\dd\mathcal A_k\|^2
    -\langle\bar\Psi_k,D_k(\mathcal A_k)\Psi_k\rangle
    -m_k\langle\bar\psi_k,\psi_k\rangle\\
   &-\varepsilon_k\Vol
    +\sum_X E_k(X,\mathcal A_k,\psi_k^\#)\Bigr\}.
 \end{aligned}
 \label{eq:digest-D-density}
\end{equation}
Here $\mathcal N_k$ abbreviates its field-independent factors
$N_kZ_kZ_k(0)$. The inserted fields are
\[
 \psi_k=H_k(\mathcal A_k)\Psi_k,\qquad
 \psi_k^\#=(\psi_k,\delta_{\alpha,\mathcal A_k}\psi_k),
\]
including barred variables. The sign of the activity in the exponent is positive,
so it contributes with a negative sign to an action defined by
$\rho_k\propto e^{-S_{\mathrm{eff},k}}$.

An essential feature of \eqref{eq:digest-D-density} is that the running
mass correction multiplies the \emph{inserted fine field}
$\bar\psi_k\psi_k$. It is not simply added to the coarse matrix
$D_k$ as a multiple of the identity. Both analyses use this
separation between a fixed-reference-mass Gaussian and a running
mass interaction.

\subsection{The manuscript's density and external slot}

\Q, Section 13.3 defines a matter density relative to the photon
Gaussian. In abbreviated notation it is
\begin{align}
 \mathscr F_k^x(A,C;\Psi)
 &=\det D_k(0)^{-1}\exp\Bigl\{
      -\langle\bar\Psi,D_k(\mathscr A+C)\Psi\rangle
      -m_k\langle\bar\psi,\psi\rangle
      -\varepsilon_k^{\mathrm v}n_k\notag\\
 &\hspace{23mm}
      +\sum_X\mathscr E_k(X;\mathscr A,C;\psi^\#)
      +\sum_X\mathfrak P_k(X,C)\Bigr\},\qquad
      \mathscr A=H_k^{\mathrm{ph}}A,
 \label{eq:digest-Q-density}\\
 &\psi=H_k(\mathscr A+C)\Psi,\qquad n_k=\ell_k^3.\notag
\end{align}
The photon quadratic weight is outside this displayed matter
density. It must be restored, or included in the reference Gaussian
measure, before comparing with \eqref{eq:digest-D-density}.
The factor $\det D_k(0)^{-1}$ is a free normalization. Its placement
differs from Dimock's accumulated normalization factors, although
both schemes track determinant ratios.

The external connection $C$ is carried in a second analytic slot;
only its sum with $\mathscr A$ enters the Dirac and field-insertion
operators. The activity and a separate scalar family retain more
information than dependence on that sum alone. The local fluctuation
determinant is split as
\begin{align}
 \delta E_k^{\mathrm{det}}(X;\mathscr A,C)
   &=E_k^d(X,\mathscr A+C)-E_k^d(X,C),\\
 E_k^{\mathrm{det}}(X,C)&=E_k^d(X,C)-E_k^d(X,0).
 \label{eq:digest-source-split}
\end{align}
The first term enters the new activity and the second accumulates
in $\mathfrak P_k$ (\Q, Section 14.7). At $C=0$,
$\mathfrak P_k(X,0)=0$ and $\mathscr E_k|_{C=0}=E_k$.
An external-field-dependent scalar cannot be discarded as a vacuum
constant when computing response derivatives. This bookkeeping is
an explicit extension of the density in \UR, (397).

Both activity constructions use balanced Grassmann kernels,
covariant H\"older-difference legs, exponential polymer decay, and
weights
\begin{equation}
 h_{k,1}=e_k^{-1/4},\qquad h_{k,2}=e_k^{-1/4+\varepsilon}.
\end{equation}
They are not literally the same Banach spaces. In \Q, Section 12.1,
the kernels have chamber labels, cluster-normalized measures,
specified admissible pair paths and boundary legs, and two connection
arguments. Its displayed norm takes a supremum of the whole weighted
coefficient norm over connections. \UR, (289)--(296), distinguishes
this type of supremum from the stronger sum of coefficientwise
suprema used in its basic norm. Bounds or identities between these
spaces require the appropriate comparison; matching the weights
alone is insufficient.

\subsection{Exactness, normalization, and the full large-field density}

The identity
$\mathbf T_k^x\mathscr F_k^x=\mathscr F_{k+1}^x$ in \Q, Section 14.7
is exact for its \emph{defined truncated transformation}, whose
photon integral is over $\mathrm{Box}_k$. It is not an identity between
an unrestricted integral and a truncated integral. Likewise, \UR's
main theorem is about the modified step (394). Recovering the full
partition function or source response requires their additional
complement arguments.

In \US\ the full density is instead a sum of terms of the form
\begin{equation}
 \rho_k
 =\sum_\Pi\mathcal N_{k,\Omega}\int\dd m_{k,\Pi}\,
 C_{k,\Pi}\chi_k(\Lambda_k)
 \exp\left\{-\tfrac12\|\dd\mathcal A_{k,\Omega}\|^2
             -S^+_{k,\Omega}+E_k(\Lambda_k)+B_{k,\Pi}\right\}.
 \label{eq:digest-history-density}
\end{equation}
Here $\dd m_{k,\Pi}$ abbreviates the retained gauge, gauge-fluctuation,
fermion, and fermion-fluctuation integrations of \US, (6).
The mass and vacuum counterterms in $S^+_{k,\Omega}$ are localized
in $\Lambda_k$; $B_{k,\Pi}$ contains boundary corrections.
There is no counterpart of all these data in the global density
\eqref{eq:digest-Q-density}. Suppressing the history and boundary
terms would erase the part of \US\ which restores the large fields.

There is a further convention issue even in Dimock's representations:
\UR's theorem is stated up to a coarse-fermion phase change
(Remark 2 after Theorem 1), and \US, Section 2.1 explicitly permits
modifications preserving the fermion integral. Equal partition
functions, equal densities after such variable changes, and equal
coefficientwise formulas are distinct claims. A literal comparison
must also match these choices.

Finally, \Q's original normalized finite-volume law has no vacuum
counterterm in its definition. A field-independent vacuum factor
cancels from normalized expectations. The tuned vacuum nevertheless
appears in the RG normalization and common pressure used to compare
partition functions on the two reference volumes (Sections 16.4,
20.4, and 22.6). This explains why the main theorem advertises a mass
assignment while its RG analysis still tunes both relevant parameters.

\section{What can, and cannot, be imported from Dimock}

The closest identification is the photon block algebra of \AG\ and
\UR\ with \Q, Section 10, at matched mesh, torus, and coarse data.
The small-field renormalization architecture in \UR\ also provides
the clear template for \Q: local propagators, determinant expansions,
connected polymer logarithms, vacuum/mass extraction, adjustment,
contraction, and a backwards/forwards fixed-point flow.

Transferring the actual estimates requires the changes to the model
and domains to be addressed. The manuscript supplies separate
sections for the chamber heavy/light decomposition, the first
fermion block integration, external-slot locality and analyticity,
and agreement on covers. These are precisely the places where a
reader should not substitute the corresponding Wilson-fermion
theorem without further argument.

The largest new dependencies are outside the common small-field RG:
the all-flux reflection-positive construction; comparison with
diagonally adapted tori; Gaussian domination in the zero-flux,
harmonic-averaged law; reference-volume control and its transfer;
and the continuum covariance and componentwise non-Gaussianity
arguments. The three supplied Dimock papers do not establish these
claims for the chamber model. Conversely, the absence of Dimock's
large-field history sum from \Q\ is not by itself a missing step if
the manuscript's proposed domination and transfer arguments hold.
Those arguments must be assessed on their own terms.

\section{Source locations for checking the comparison}

The following labels are literal labels in \texttt{QED3/qed3.tex}.
They allow the most consequential claims above to be located without
depending on the equation numbering of a separately compiled PDF.

\small
\begin{longtable}{@{}P{0.29\textwidth}P{0.65\textwidth}@{}}
\toprule
Topic and manuscript section & Searchable label\\
\midrule
\endfirsthead
\toprule Topic and manuscript section & Searchable label\\\midrule
\endhead
Main conclusion, \S1 & \qlabel{thm:main}\\[4pt]
Flux quotient, \S2.1 & \qlabel{eq:reg-configuration}\\[4pt]
Chamber operator, \S2.3 & \qlabel{eq:reg-dirac}\\[4pt]
Light Schur complement, \S9.2 & \qlabel{prop:schur-symbol}\\[4pt]
Photon block decomposition, \S10.5 &
\qlabel{eq:photon-Gaussian-block-decomposition}\\[4pt]
First fermion integration, \S11.5 & \qlabel{eq:Gaussian-block-integration}\\[4pt]
Relevant extraction, \S12.3 & \qlabel{eq:rg-exact-extraction-identity}\\[4pt]
Normalized matter density, \S13.3 & \qlabel{eq:rg-step-input-density}\\[4pt]
Truncated block map, \S13.3 & \qlabel{eq:rg-step-block-transformation}\\[4pt]
Nonlinear step, \S14.7 & \qlabel{thm:rgnl-current-step}\\[4pt]
Common bare pair, \S16.4 & \qlabel{thm:rgfl-common-pair}\\[4pt]
Zero-sector unfolding, \S20.3 & \qlabel{eq:external-zero-sector}\\[4pt]
Zero-sector domination, \S21.5 &
\qlabel{prop:interacting-zero-flux-domination}\\[4pt]
Restricted terminal density, \S22.2 & \qlabel{thm:term-density}\\[4pt]
Complement estimate, \S22.3 & \qlabel{prop:term-complement}\\[4pt]
Common pressure, \S22.6 & \qlabel{thm:term-common-partition}\\[4pt]
Response complement, \S23.4 & \qlabel{prop:interacting-complement}\\[4pt]
Volume transfer, \S25.1 & \qlabel{prop:transfer-general}\\[4pt]
Flux contribution, \S25.3 & \qlabel{prop:transfer-flux-defect}\\
\bottomrule
\end{longtable}
\normalsize

For Dimock, the main checkpoints are \AG, Sections 2.4, 3.3--3.5,
and 3.10; \UR, Sections 1.2--1.4, 2.2--2.5, 3.1--3.2, 4.2--4.6,
6.1, and 7; and \US, Sections 2.1--2.3 and 4.1--4.3.
The supplied PDFs carry arXiv version stamps v3, v3, and v2,
respectively. The comparisons above use those files, including their
equation numbers, rather than assuming that the first submission
dates in their filenames identify the versions.

\begin{thebibliography}{Q}
\raggedright
\bibitem[Q]{Q}
\emph{The field-strength Osterwalder--Schrader axioms for QED in three
dimensions}, local manuscript \path{QED3/qed3.tex}, as present on
22 October 2026.

\bibitem[A]{A}
J.~Dimock, \emph{Covariant Axial Gauge},
\href{https://arxiv.org/abs/1410.2798v3}{arXiv:1410.2798v3}.
Supplied file: \path{Dimock/Dimock14102798CovariantAxialGauge.pdf}.

\bibitem[R]{R}
J.~Dimock, \emph{Ultraviolet regularity for QED in $d=3$},
\href{https://arxiv.org/abs/1512.04373v3}{arXiv:1512.04373v3}.
Supplied file: \path{Dimock/Dimock151204373UVRegularityforQED.pdf}.

\bibitem[S]{S}
J.~Dimock, \emph{Ultraviolet stability for QED in $d=3$},
\href{https://arxiv.org/abs/2009.01156v2}{arXiv:2009.01156v2}.
Supplied file: \path{Dimock/Dimock200901156UVStabilityforQED.pdf}.
\end{thebibliography}
\end{document}
